\documentclass[12pt]{article}
\usepackage[T1]{fontenc}
\usepackage[utf8]{inputenc}
\usepackage[polish]{babel}
\usepackage{lmodern}
\usepackage[a4paper,margin=2.4cm]{geometry}
\usepackage{xcolor}
\usepackage{graphicx}
\usepackage{amsmath,amssymb}
\usepackage{booktabs}
\usepackage{array}
\usepackage{enumitem}
\usepackage{hyperref}
\usepackage{tcolorbox}
\usepackage{fancyhdr}
\definecolor{WEIIPrimary}{RGB}{0,130,150}
\definecolor{WEIIPrimaryDark}{RGB}{0,86,100}
\definecolor{WEIIPrimaryLight}{RGB}{229,246,248}
\definecolor{PRzNavy}{RGB}{31,63,142}
\definecolor{RLTextDark}{RGB}{34,34,34}
\definecolor{RLGray}{RGB}{90,90,90}
\definecolor{RLLightGray}{RGB}{245,246,248}
\definecolor{RLWarning}{RGB}{184,86,0}
\definecolor{RLExample}{RGB}{52,112,64}
\definecolor{RLCodeBg}{RGB}{248,248,248}

\usepackage{listings}
\usepackage{xcolor}

\lstdefinestyle{rlpython}{%
    language=Python,
    basicstyle=\ttfamily\scriptsize,
    keywordstyle=\color{WEIIPrimaryDark}\bfseries,
    commentstyle=\color{RLGray}\itshape,
    stringstyle=\color{RLExample},
    numbers=left,
    numberstyle=\tiny\color{RLGray},
    numbersep=6pt,
    showstringspaces=false,
    breaklines=true,
    tabsize=4,
    frame=single,
    rulecolor=\color{WEIIPrimary!35},
    backgroundcolor=\color{RLCodeBg},
    captionpos=b
}

\lstset{style=rlpython}


\hypersetup{colorlinks=true,linkcolor=WEIIPrimaryDark,urlcolor=WEIIPrimaryDark,citecolor=WEIIPrimaryDark}
\setlength{\headheight}{15pt}
\pagestyle{fancy}
\fancyhf{}
\lhead{Uczenie się ze wzmocnieniem}
\rhead{WEiI PRz}
\cfoot{\thepage}

\newtcolorbox{taskbox}[1]{colback=WEIIPrimaryLight,colframe=WEIIPrimary,title=\textbf{#1}}
\newtcolorbox{notebox}[1]{colback=RLLightGray,colframe=RLGray,title=\textbf{#1}}
\newtcolorbox{warninghandout}[1]{colback=RLWarning!8,colframe=RLWarning,title=\textbf{#1}}

\newcommand{\handouttitle}[3]{%
\begin{center}
{\Large\bfseries #1}\\[0.2cm]
{\large #2}\\[0.15cm]
{\small #3}
\end{center}
\vspace{0.4cm}
\hrule
\vspace{0.6cm}
}

\newcommand{\EnableLabNumbering}{%
  \renewcommand{\thesection}{L\arabic{section}}%
  \renewcommand{\thesubsection}{\thesection.\arabic{subsection}}%
  \renewcommand{\thesubsubsection}{\thesubsection.\arabic{subsubsection}}%
}


\title{Raport projektowy\\Uczenie się ze wzmocnieniem}
\author{Patryk Dziki}
\date{2026/2027}

\begin{document}
\maketitle
\tableofcontents
\newpage

\section{Informacje podstawowe}

\begin{center}
\begin{tabular}{p{0.34\textwidth}p{0.56\textwidth}}
\toprule
Element & Informacja \\
\midrule
Imię i nazwisko & Patryk Dziki \\
Numer albumu & 179480 \\
Kierunek/specjalność & Informatyka, specjalność: AI \\
Tytuł projektu & Sterowanie wysokością drona z ograniczoną energią z~wykorzystaniem uczenia się ze wzmocnieniem \\
Link do repozytorium & \url{https://github.com/vvking321/rl-drone-control} \\
Data oddania & Do uzupełnienia \\
\bottomrule
\end{tabular}
\end{center}

\section{Temat i cel projektu}

\subsection{Opis problemu}

Projekt dotyczy sterowania pionowym ruchem drona (w jednym wymiarze). Dron ma wznieść się na zadaną wysokość i ją utrzymać, zużywając przy tym jak najmniej energii. Jest on modelowany jako punkt materialny poruszający się w pionie pod wpływem ciągu silników, grawitacji i losowych podmuchów wiatru, a jego bateria jest ograniczona. W kolejnych krokach czasowych sterownik wybiera, czy poziom ciągu ma zostać zmniejszony, utrzymany czy zwiększony.

Problem jest uproszczonym modelem zadania występującego w rzeczywistych dronach. Pozwala przebadać proces uczenia agenta, wpływ konstrukcji funkcji nagrody oraz porównać agenta uczącego się z klasycznym regulatorem.

\subsection{Cel agenta}

Agent ma doprowadzić drona na zadaną wysokość i utrzymać go w strefie docelowej, nie doprowadzając do rozbicia, wyjścia poza dopuszczalny zakres wysokości ani rozładowania baterii, i zużywając przy tym jak najmniej energii.

Epizod uznaje się za udany, jeśli zakończył się bez porażki, a dron przebywał w strefie docelowej przez co najmniej 80\% kroków (wartość wstępna).

Sukces projektu będzie oceniany na podstawie: średniego zwrotu z epizodu, odsetka epizodów udanych, zużycia energii, liczby porażek oraz porównania z agentem bazowym (regulator PID i prosta reguła progowa).

\subsection{Uzasadnienie zastosowania uczenia się ze wzmocnieniem}

Problem można rozpatrywać jako problem uczenia się ze wzmocnieniem z następujących powodów:
\begin{itemize}
    \item \textbf{Decyzje sekwencyjne:} w każdym kroku czasowym agent wybiera zmianę ciągu, a całe zadanie jest ciągiem takich decyzji.
    \item \textbf{Wpływ na przyszły stan:} zmiana ciągu wpływa na prędkość, a dopiero potem na wysokość (bezwładność), a ponadto zużywa baterię, co ogranicza możliwości w kolejnych krokach.
    \item \textbf{Wynik długoterminowy:} agresywne zwiększenie ciągu szybko zbliża drona do celu, ale może prowadzić do przeregulowania i dużego zużycia energii. Dobra strategia musi uwzględniać skumulowany koszt całego epizodu, a nie tylko natychmiastową nagrodę.
    \item \textbf{Zasadność agenta uczącego się:} agent uczy się strategii wyłącznie z interakcji ze środowiskiem, bez znajomości modelu dynamiki i bez ręcznego strojenia, a jego kompromis między dokładnością a energią wynika z funkcji nagrody. Zakłócenia (wiatr) nie są obserwowane, więc strategia musi działać dobrze w ich obecności.
\end{itemize}

\section{Formalizacja problemu RL/MDP}

\subsection{Agent i środowisko}

\textbf{Agent} jest sterownikiem ciągu drona. \textbf{Środowisko} stanowi jednowymiarowy model dynamiki pionowej z zakłóceniem wiatru i modelem zużycia baterii. Dron startuje z losowej wysokości początkowej $z_0 > 0$ (z przedziału, którego granice zostaną ustalone przy implementacji) i zerowej prędkości, a docelowa wysokość wynosi $z_{\mathrm{ref}}$.

Jeden krok epizodu przebiega następująco:
\begin{enumerate}
    \item agent obserwuje stan i wybiera akcję $a_t \in \{0, 1, 2\}$,
    \item środowisko aktualizuje poziom ciągu,
    \item losowane jest zakłócenie $w_t$, a następnie aktualizowane są prędkość i wysokość,
    \item aktualizowany jest poziom baterii,
    \item obliczana jest nagroda $r_t$ i sprawdzane są warunki zakończenia epizodu,
    \item środowisko zwraca nową obserwację, nagrodę oraz informację o zakończeniu epizodu.
\end{enumerate}

Dynamika środowiska jest opisana słownie, a dokładne wzory zostaną podane przy implementacji środowiska. W każdym kroku czasowym środowisko zmienia poziom ciągu zgodnie z wybraną akcją, a z ciągu, grawitacji i losowego zakłócenia (podmuchu wiatru) wylicza nową prędkość, a z niej nową wysokość. Poziom baterii maleje proporcjonalnie do użytego ciągu. Błąd wysokości definiujemy jako $e_t = z_t - z_{\mathrm{ref}}$, a $e_{\max}$ oznacza maksymalny dopuszczalny błąd wysokości (wstępnie 5~m).

\subsection{Stan lub obserwacja}

Pełny stan środowiska obejmuje wysokość, prędkość pionową, ciąg i poziom baterii oraz bieżące zakłócenie $w_t$. Agent \textbf{nie obserwuje zakłócenia $w_t$}, więc jego obserwacja jest niepełna (różni się od stanu środowiska). W wersji tablicowej każda ciągła składowa obserwacji zostaje zdyskretyzowana na skończoną liczbę przedziałów, a rozdzielczość dyskretyzacji będzie przedmiotem eksperymentu.

\begin{center}
\begin{tabular}{p{0.30\textwidth}p{0.60\textwidth}}
\toprule
Składowa stanu/obserwacji & Znaczenie \\
\midrule
Błąd wysokości $e_t$ [m] & Różnica między wysokością drona a wysokością docelową ($e_t = z_t - z_{\mathrm{ref}}$); wartość dodatnia oznacza położenie powyżej celu. \\
Prędkość pionowa $v_t$ [m/s] & Prędkość drona w pionie; wartość dodatnia oznacza ruch w górę. \\
Poziom baterii $b_t$ & Pozostała energia, znormalizowana do przedziału $[0,1]$. \\
Poziom ciągu $u_t$ & Aktualny, znormalizowany ciąg silników, $T_t = u_t\,T_{\max}$. \\
\bottomrule
\end{tabular}
\end{center}

\subsection{Akcje}

Agent wybiera jedną z trzech akcji dyskretnych, które zmieniają poziom ciągu o stały krok $\delta_u$. Żadna akcja nie jest zabroniona. Jeżeli akcja miałaby wyprowadzić ciąg poza przedział $[0,1]$, poziom ciągu zostaje obcięty do granicy (akcja nie ma wtedy dodatkowego efektu).

\begin{center}
\begin{tabular}{p{0.30\textwidth}p{0.60\textwidth}}
\toprule
Akcja & Opis \\
\midrule
$a=0$: zmniejsz ciąg & $u_{t+1} = \max(0,\ u_t - \delta_u)$ \\
$a=1$: utrzymaj ciąg & $u_{t+1} = u_t$ \\
$a=2$: zwiększ ciąg & $u_{t+1} = \min(1,\ u_t + \delta_u)$ \\
\bottomrule
\end{tabular}
\end{center}

\subsection{Funkcja nagrody}

Nagroda w kroku $t$ jest sumą czterech składników:
\begin{equation}
    r_t = r_t^{\text{błąd}} + r_t^{\text{energia}} + r_t^{\text{cel}} + r_t^{\text{porażka}},
\end{equation}
gdzie:
\begin{itemize}
    \item $r_t^{\text{błąd}} = -\alpha_e\,|e_t| / e_{\max}$ --- kara proporcjonalna do odchylenia od wysokości docelowej,
    \item $r_t^{\text{energia}} = -\alpha_u\,u_t$ --- kara proporcjonalna do użytego ciągu (koszt energii),
    \item $r_t^{\text{cel}} = r_{\mathrm{cel}}$, jeśli $|e_t| \le \varepsilon_e$ i $|v_t| \le \varepsilon_v$, a w przeciwnym razie $0$ --- bonus za utrzymanie wysokości,
    \item $r_t^{\text{porażka}} = -R_{\mathrm{fail}}$ w kroku, w którym epizod kończy się porażką, a w przeciwnym razie $0$.
\end{itemize}
W definicji nagrody wielkości $e_t$, $v_t$ i $u_t$ odnoszą się do stanu po wykonaniu kroku (kolejność operacji opisano w podrozdziale „Agent i środowisko"). Wstępne wartości parametrów: $\alpha_e = 1$, $\alpha_u = 0{,}1$, $r_{\mathrm{cel}} = 1$, $R_{\mathrm{fail}} = 100$, $\varepsilon_e = 0{,}1$~m, $\varepsilon_v = 0{,}2$~m/s. Zostaną one skorygowane po pierwszych testach.

\subsection{Warunki zakończenia epizodu}

Epizod kończy się (dla stanu po wykonaniu kroku), gdy:
\begin{itemize}
    \item $z_t \le 0$ (rozbicie drona o ziemię) --- porażka,
    \item $|e_t| > e_{\max}$ (wyjście poza dopuszczalny zakres wysokości) --- porażka,
    \item $b_t \le 0$ (wyczerpanie baterii) --- porażka,
    \item osiągnięto limit $N_{\max}$ kroków (wstępnie 400, czyli 20~s przy $\Delta t = 0{,}05$~s) --- koniec neutralny.
\end{itemize}

\subsection{Założenia i ograniczenia środowiska}

Przyjęto następujące uproszczenia:
\begin{itemize}
    \item ruch tylko w pionie (jeden wymiar), bez orientacji, obrotów i ruchu poziomego,
    \item dron jako punkt materialny o stałej masie, bez aerodynamiki i efektu przyziemienia,
    \item ciąg zmienia się natychmiast po wyborze akcji (brak opóźnienia silników),
    \item prosty model wiatru jako losowego zakłócenia przyspieszenia oraz liniowy model zużycia baterii,
    \item akcje dyskretne o stałym kroku zmiany ciągu.
\end{itemize}

Model nie odwzorowuje rzeczywistego drona, więc wyniki nie dowodzą skuteczności sterowania na prawdziwym sprzęcie.

Elementy formalizacji, które mogą jeszcze wymagać korekty po pierwszych testach: wartości parametrów fizycznych, składniki i wagi funkcji nagrody, rozdzielczość dyskretyzacji obserwacji, krok zmiany ciągu $\delta_u$, definicja udanego epizodu oraz limit kroków.

\section{Implementacja}

\subsection{Struktura repozytorium}

Należy krótko opisać strukturę repozytorium i najważniejsze pliki projektu.

\begin{verbatim}
repozytorium/
|-- src/
|-- experiments/
|-- report/
|-- results/
`-- README.md
\end{verbatim}

\subsection{Środowisko}

Należy opisać sposób implementacji środowiska, w szczególności funkcje inicjalizacji epizodu oraz wykonania kroku środowiska.

\subsection{Agent bazowy}

Należy opisać agenta losowego, heurystycznego albo inny baseline wykorzystany do porównania z agentem uczącym się.

\subsection{Agent uczący się}

Należy opisać zastosowany algorytm RL lub DRL, jego parametry oraz sposób implementacji. W przypadku Q-tab należy opisać reprezentację tablicy wartości. W przypadku DRL należy opisać wejścia modelu, wyjścia modelu, architekturę oraz sposób uczenia.

\section{Plan eksperymentów}

Należy opisać eksperymenty zaplanowane w celu oceny działania agenta.

\begin{center}
\begin{tabular}{p{0.24\textwidth}p{0.30\textwidth}p{0.30\textwidth}}
\toprule
Eksperyment & Co jest zmieniane? & Co jest mierzone? \\
\midrule
E1 & Do uzupełnienia & Do uzupełnienia \\
E2 & Do uzupełnienia & Do uzupełnienia \\
E3 & Do uzupełnienia & Do uzupełnienia \\
\bottomrule
\end{tabular}
\end{center}

W tej sekcji należy wskazać:
\begin{itemize}
    \item liczbę epizodów uczenia,
    \item liczbę powtórzeń eksperymentu albo uzasadnienie braku powtórzeń,
    \item parametry uczenia,
    \item metryki oceny,
    \item warianty środowiska, map lub funkcji nagrody,
    \item sposób porównania z agentem bazowym.
\end{itemize}

\section{Wyniki}

W tej sekcji należy przedstawić wyniki eksperymentów. Wyniki powinny być pokazane w postaci tabel, wykresów lub czytelnych zestawień.

\begin{center}
\begin{tabular}{p{0.28\textwidth}p{0.20\textwidth}p{0.20\textwidth}p{0.20\textwidth}}
\toprule
Wariant & Średni zwrot & Skuteczność & Długość epizodu \\
\midrule
Agent bazowy & Do uzupełnienia & Do uzupełnienia & Do uzupełnienia \\
Agent RL/DRL & Do uzupełnienia & Do uzupełnienia & Do uzupełnienia \\
\bottomrule
\end{tabular}
\end{center}

\begin{figure}[h]
    \centering
    \fbox{\parbox{0.82\textwidth}{\centering Miejsce na wykres procesu uczenia}}
    \caption{Przykładowy wykres procesu uczenia agenta.}
\end{figure}

\section{Analiza i wnioski}

Należy zinterpretować wyniki. W tej części należy odpowiedzieć m.in. na pytania:
\begin{itemize}
    \item Czy agent uczący się działa lepiej niż agent bazowy?
    \item Jak zmieniało się zachowanie agenta w trakcie uczenia?
    \item Które elementy środowiska były najtrudniejsze?
    \item Czy funkcja nagrody prowadziła do oczekiwanego zachowania?
    \item Czy wystąpiły przypadki niepożądanego zachowania agenta?
    \item Jakie są najważniejsze ograniczenia rozwiązania?
\end{itemize}

\section{Instrukcja uruchomienia}

Należy opisać sposób uruchomienia projektu od czystego repozytorium. Przykładowo:

\begin{verbatim}
git clone https://github.com/...
cd nazwa-repozytorium
python -m venv .venv
.venv\Scripts\activate
pip install -r requirements.txt
python src/main.py
\end{verbatim}

Należy również wskazać, jak odtworzyć eksperymenty oraz gdzie zapisywane są wyniki.

\section{Wsparcie AI}

W tej sekcji należy obowiązkowo opisać wykorzystanie narzędzi sztucznej inteligencji. Jeżeli student nie korzystał z takich narzędzi, należy wpisać jednoznaczną informację, że nie korzystano ze wsparcia AI.

\begin{center}
\begin{tabular}{p{0.13\textwidth}p{0.19\textwidth}p{0.21\textwidth}p{0.27\textwidth}}
\toprule
Data & Narzędzie/model & Cel użycia & Transkrypt / link / zakres \\
\midrule
Do uzupełnienia & Do uzupełnienia & Do uzupełnienia & Do uzupełnienia \\
Do uzupełnienia & Do uzupełnienia & Do uzupełnienia & Do uzupełnienia \\
\bottomrule
\end{tabular}
\end{center}

Należy zapisać transkrypty rozmów z modelami AI wraz z datą oraz nazwą modelu. Jeżeli rozmowa odbywała się poza narzędziami programistycznymi, należy dodatkowo udostępnić całą rozmowę poprzez link, o ile użyte narzędzie umożliwia takie udostępnienie.

Student powinien także wyjaśnić, które elementy wygenerowane lub zasugerowane przez AI zostały samodzielnie sprawdzone, zmodyfikowane albo odrzucone.

\section{Literatura}

W tej sekcji należy wskazać wykorzystane źródła: książki, artykuły naukowe, dokumentację bibliotek, materiały wykładowe oraz inne źródła użyte podczas pracy nad projektem.

\begin{thebibliography}{9}
\bibitem{sutton_barto}
R. S. Sutton, A. G. Barto, \textit{Reinforcement Learning: An Introduction}, 2nd ed., MIT Press, 2018.
\end{thebibliography}

\end{document}
